\scanpage{11}
\mspage{10}
\mspage{IV-\unsure{964|764}}

\uline{Pour que} $A$ \uline{soit unibranche} (resp. \uline{géométriquement unibranche), il faut et il suffit que} $\operatorname{Spec}(A)$ \uline{soit irréductible}.

\uline{Proposition} (18.8.1\deleted{3}\inserted{5}).-- \uline{Soit} $A$ \uline{un anneau local noethérien. Si} $A$ \uline{est universellement caténaire, il en est de même de} ${}^{\mathrm{hs}}A$.

La démonstration est la même que celle de (18.7.5).

\uline{18.9. Fibres formelles des anneaux noethériens henséliens}.

\uline{Théorème} (18.9.1).-- \uline{Soit} $A$ \uline{un anneau local noethérien hensélien, dont les fibres formelles sont géométriquement normales. Alors les fibres formelles de} $A$ \uline{sont géométriquement intègres} (donc \uline{géométriquement connexes}).

Les deux assertions de l'énoncé sont en fait équivalentes, un préschéma \inserted{localement noethérien} connexe et normal étant intègre. Comme toute $A$-algèbre intègre finie est un anneau local hensélien (18.5.9 et 18.5.10), il résulte de (7.3.16.2) que l'on est ramené à prouver que si l'on suppose de plus $A$ \uline{intègre}, alors la fibre du morphisme $\operatorname{Spec}(\hat A)\to\operatorname{Spec}(A)$ au point générique de $\operatorname{Spec}(A)$ est \uline{intègre}, ce qui résulte du

\uline{Corollaire} (18.9.2).-- \uline{Soit} $A$ \uline{un anneau local noethérien hensélien intègre, dont les fibres formelles sont géométriquement normales. Alors} $\hat A$ \uline{est inté-}\deleted{\illegible{1}}\uline{gre}.

En effet, il résulte de (18.8.12) que $A$ est unibranche, donc \deleted{\illegible{1}} le théorème résulte de l'hypothèse faite sur les fibres formelles de $A$ et de (7.6.3).

\uline{Corollaire} (18.\deleted{\unsure{8}}\inserted{9}.3).-- \uline{Sous les hypothèses de} (18.9.2), \uline{le corps des fractions} $L$ \uline{de} $\hat A$ \uline{est une extension séparable du corps des fractions} $K$ \uline{de} $A$, \uline{et} $K$ \uline{est algébriquement fermé dans} $L$.

\deleted{\unsure{La première assertion}} \deleted{\illegible{4}} \inserted{Cela} résulte de ce que la fibre du morphisme $\operatorname{Spec}(\hat A)\to{}\to\operatorname{Spec}(A)$ au point générique de $\operatorname{Spec}(A)$ est géométriquement \deleted{\unsure{normal}}intègre et de (4.3.2) \inserted{et (4.3.5)}.

